\documentclass[oneside]{scrbook}
\KOMAoptions{
  paper=letter,
  fontsize=12pt,
  parskip=full,
  headings=small,
}

\usepackage{import}
\usepackage{tikz}

\usetikzlibrary{arrows.meta, positioning, calc, decorations.pathreplacing}

\usepackage[mode=tex]{standalone}

\usepackage{amsmath, amssymb, amsthm, mathtools, stmaryrd}
\usepackage{algorithm, algorithmic}

\usepackage{booktabs}

\newtheorem{theorem}{Theorem}[chapter]
\newtheorem{proposition}[theorem]{Proposition}
\newtheorem{lemma}[theorem]{Lemma}
\newtheorem{corollary}[theorem]{Corollary}
\theoremstyle{definition}
\newtheorem{definition}[theorem]{Definition}
\newtheorem{example}[theorem]{Example}
\newtheorem{remark}[theorem]{Remark}

\usepackage{fmtcount}

\usepackage{mlmodern}
\usepackage[T1]{fontenc}
\usepackage{microtype}

\usepackage{pgffor}

\usepackage{xcolor}
\usepackage[colorlinks, allcolors=blue!50!black]{hyperref}

\usepackage{cleveref}

\usepackage{enumitem}
\setlist{nosep}

\pagestyle{plain}

\usepackage[backend=biber, style=alphabetic]{biblatex}
\addbibresource{references.bib}

\usepackage[en-CA]{datetime2}

\newcommand{\lecture}[3]{
    \clearpage
    \phantomsection
    \noindent
    \Ordinalstringnum{#1} Lecture \hfill \DTMdate{#2}\\[0.5em]
    {\Large \textsf{\textbf{#3}}}
    \par\vspace{0.5em}
    \hrule
    \vspace{1em}
    \addcontentsline{toc}{chapter}{\texorpdfstring{\Ordinalstringnum{#1} Lecture}{Lecture #1}: #3 (\texorpdfstring{\DTMdate{#2}}{#2})}
    \stepcounter{chapter}
    \setcounter{chapter}{#1}
}

\title{Notes for Public-Key Cryptography}
\author{Guilherme Zeus Dantas e Moura}
\date{Fall 2026}

\begin{document}

\maketitle

This document contains my lecture notes for CO 685,
a graduate course on The Mathematics of Public-Key Cryptography,
taught by Sam Jaques at the University of Waterloo in the Fall 2026 term.
Any errors, inaccuracies, or misinterpretations are my own.
Use these notes at your own discretion,
and please feel free to point out any corrections that may be needed.

Here's the course description from the outline, by Sam Jaques:
\begin{quote}
    An in-depth study of public-key cryptography, including: number-theoretic problems - prime generation, integer factorization, discrete logarithms; public-key encryption; digital signatures; key establishment; elliptic curve cryptography; post-quantum cryptography; and security proofs.
\end{quote}

The tentative class plan is to cover:
\begin{itemize}[left=0pt]
    \item Diffie-Hellman Key Exchange (\(\approx 6\) weeks)
    \begin{itemize}[left=0pt]
        \item The basic idea of public key exchange over groups; the discrete log problem; square-and-multiply
        \item The group of \(\mathbb{Z}_p^*\): primality testing; structure of modular integers; rings and fields
        \item Elliptic curves: a geometric picture; unexplained facts about elliptic curves
        \item Discrete log algorithms: Baby-step giant-step, Pollard's rho, Pohlig-Hellman, Index Calculus
    \end{itemize}
    \item Formal Security (\(\approx 4\) weeks)
    \begin{itemize}[left=0pt]
        \item Working with security assumptions; assumptions over finite fields and elliptic curves
        \item IND-CPA and IND-CCA; proofs of IND-C\{C/P\}A security of Diffie-Hellman exchange; Fujisaki-Okamoto transform
        \item Zero Knowledge proofs; Fiat-Shamir transform
        \item Digital signatures (Schnorr, DSA, ECDSA)
    \end{itemize}
    \item Learning-With-Errors (\(\approx 2\) weeks)
    \begin{itemize}[left=0pt]
        \item LWE definition; key exchange; signatures
        \item Lattices; lattice attacks
        \item Ring-LWE; polynomial rings and finite fields
    \end{itemize}
    \item RSA (time permitting)
    \begin{itemize}[left=0pt]
        \item Encryption; signatures
        \item Attacks; security proofs
    \end{itemize}
\end{itemize}

The references for the course are the following textbooks and resources:
\cite{HoffsteinPipherSilverman2014},
\cite{BonehShoup2023},
\cite{JaquesNotes}.

\tableofcontents

\foreach \i in {01, 02, 03, 04, 05, 06, 07, 08, 09, 10, ..., 99} {%
    \edef\FileName{lectures/\i.tex}%
    \IfFileExists{\FileName}{%
      \input{\FileName}%
    }%
  }

\printbibliography

\end{document}
